\subsection{The four-loop dyadic example with unequal slopes}
In Example~(3) of Section~7 of DuBose and Valli\`eres, the four directions are $(1,5),(0,3),(1,2),(0,1)$. Write $T_n=v_2(\kappa_n)$ for its spanning-tree valuation. Then
\[
 T_n=n2^n+\tfrac{33}{4}2^n-4n-1\qquad(n\ge4),
 \qquad (T_1,T_2,T_3)=(5,19,65).
\]
To see this, write $F=8-\sum_{(u,v)}(X^uY^v+X^{-u}Y^{-v})$.
At conductor $2^s$, take $\zeta$ primitive, $e=2^{s-1}$, and let $W_s$ be the sum of $v_2(F(\zeta^a,\zeta^b))$ over pairs modulo $2^s$ with at least one odd entry. The two vectors $(1,2),(0,1)$ generate the group, so the only zero of $F$ among these characters is the excluded trivial character. The matrix-tree theorem and character diagonalization give $T_s-T_{s-1}=W_s-2$.

The decisive local computation is the following identity over $\mathbb F_2$:
\[
 XY^5\overline F=(Y+1)\big[(Y^2+Y+1)(X^2Y^7+1)+XY^2(Y+1)^5\big].
\]
Substitute $X=(1+z)^a$, $Y=(1+z)^b$, using nonnegative representatives of $a,b$. For $b\ne0$, put $B=2^{v_2(b)}$, $C=2^{v_2(2a+7b)}$. In characteristic two, $(1+z)^j-1$ has order $2^{v_2(j)}$. Since $C$ is a power of two, $C\ne5B$. The resulting order of $\overline F$ is exactly $D=B+\min(C,5B)$. Whenever $D<e$ this gives $ev_2(F(\zeta^a,\zeta^b))=D$: in the integral expansion at $z=\zeta-1$ the lower powers have even coefficients, the coefficient of $z^D$ is odd, and $v_2(\zeta-1)=1/e$. One can use the polynomial $XY^5 F$ to avoid negative exponents throughout.

For $s\ge5$, the sums over the following disjoint classes exhaust the conductor shell.
\begin{center}
\begin{tabular}{lll}
\hline
Class & Normalized order $ev_2(F)$ & Contribution to $W_s$\\
\hline
$b$ odd & $2$ & $4e$\\
$b\equiv2\pmod4$ & $6,10,12$ & $17e/4$\\
$v_2(b)=k$, $2\le k\le s-2$ & $2^k+2$ & $(s-3)e+e-4$ (all $k$)\\
$b=0$ & $e+2$ & $e+2$\\
$b=e$ & $2e$ & $2e$\\
\hline
\end{tabular}
\end{center}
Indeed, for fixed $b\equiv2\pmod4$, as $a$ runs through odd residues, $C$ is $4,8,$ or at least $16$ in proportions $1/2,1/4,1/4$. For $k\ge2$ we have $C=2$ and $e/2^k$ eligible values of $b$ at depth $k$. All three generic rows have $D<e$. For the two exceptional rows, the integer identities $F(X,1)=2(2-X-X^{-1})$ and $F(X,-1)=12$ give the listed orders directly. Summation yields
\[
 W_s=(s+37/4)2^{s-1}-2\quad(s\ge5).
\]
The small-shell totals $W_1,W_2,W_3,W_4$ are $7,16,48,116$. For an exact short verification, use one representative per Galois orbit, $(a,1)$ with arbitrary $a$ and $(1,b)$ with $b$ even. Reduce $F(Z^a,Z^b)$ modulo $Z^e+1$ and write the result as $\sum_{j<e}c_j(Z-1)^j$. The orbit contributes $\min_{c_j\ne0}(ev_2(c_j)+j)$, since its size is $e$ and these orders are distinct modulo $e$. At $s=1,2,3,4$, the multisets of orbit contributions, using superscripts for multiplicities, are respectively
\[
 \{2^{[2]},3\},\quad \{2^{[4]},4^{[2]}\},\quad
 \{2^{[8]},6^{[2]},8,12\},\quad
 \{2^{[16]},6^{[4]},10^{[2]},16,24\}.
\]
Thus $T_4=179$. Summing the shell recurrence proves the formula for all subsequent levels.
